% ECONOMICS_PROBLEM: {"id": "J23-400", "title": "AI and automation: displacement versus complementarity", "jel_code": "J23", "source_id": "OP-0022"}
% !TeX program = xelatex
\documentclass[11pt,a4paper]{article}
\usepackage[margin=23mm]{geometry}
\usepackage{fontspec}
\setmainfont{Latin Modern Roman}
\usepackage{amsmath,amssymb,mathtools}
\usepackage{xcolor,enumitem,booktabs,longtable,array,xurl,hyperref}
\definecolor{muted}{HTML}{53636C}
\hypersetup{colorlinks=true,urlcolor=blue,linkcolor=blue}
\setlength{\parindent}{0pt}
\setlength{\parskip}{5pt}
\newcommand{\R}{\mathbb R}
\newcommand{\E}{\mathbb E}
\newcommand{\N}{\mathbb N}
\newcommand{\ind}{\mathbf 1}
\newcommand{\Var}{\operatorname{Var}}
\newcommand{\Cov}{\operatorname{Cov}}
\newcommand{\supp}{\operatorname{supp}}
\DeclareMathOperator*{\argmin}{arg\,min}
\DeclareMathOperator*{\argmax}{arg\,max}
\newcommand{\entrymeta}[3]{{\sffamily\small\color{muted}\textbf{\texttt{#1}}\quad|\quad #2\par #3\par}}
\newcommand{\Problemref}[1]{\texttt{#1}}
\newcommand{\JELref}[2]{\texttt{#2} (JEL \texttt{#1})}
\begin{document}
\section*{AI and automation: displacement versus complementarity}
\textbf{Problem ID:} \texttt{J23-400}\quad \textbf{JEL:} \texttt{J23}\par
% BEGIN ECONOMICS STATEMENT
\label{problem:OP-0022}
\label{atlas:prob:L2}
\noindent\textbf{Original atlas ID:} L2 (entry 22). \textbf{Classification:} Editorial dynamic equilibrium prediction and identification problem.

\noindent\textbf{Original atlas question:} Which tasks disappear, which expand, and how rapidly do firms create new work?

\subsection*{Definitions and assumptions}
Fix a country, a time horizon $H\in\mathbb N$, and an adoption policy $a$ specifying access, prices, and permitted uses of an AI technology. A dynamic model $M$ contains workers with measured baseline types $X$, firms choosing adoption and task organization, a set of task labels $\mathcal K_t$ at each date drawn from a fixed common task universe, and laws for entry, capital accumulation, skill acquisition, and invention. Specify production functions, information, adjustment costs, market-clearing equations, and initial conditions. Assume an equilibrium exists; impose uniqueness or state an explicit selection rule. Let $L_{kt}^M(a)$ be paid human hours on task $k$, including zero for an absent task, and let $E_t^M(a)$ and $W_t^M(a)$ denote employment and the wage distribution.
\subsection*{Formal research question}
Relative to a specified baseline $a_0$, characterize the identified set of the equilibrium trajectory
\[
 \Theta_H(M)=\bigl(
 E_t^M(a)-E_t^M(a_0),\;
 W_t^M(a),\;
 (L_{kt}^M(a)-L_{kt}^M(a_0))_k,\;
 N_t^M(a)\bigr)_{t=0}^H,
\]
where $N_t^M(a)$ is the measure of newly created task labels under an announced, stable classification rule. For every baseline worker type $x$ with positive support, estimate the change in discounted labor earnings, including unemployment spells. Determine which adoption designs and equilibrium restrictions distinguish productivity-driven task expansion from substitution of machines for human tasks, allowing both to occur simultaneously.
\subsection*{Interpretation and scope}
Task exposure is a predictor, not randomized adoption. A workplace experiment with assignment $Z$ can identify $\mathbb E[Y_i(1)-Y_i(0)]$ for its participants under its interference assumptions, where $Y_i$ is measured output; that effect does not equal any component of $\Theta_H$. Firms can change hiring, prices, product quality, and entry after the experiment. New tasks require a reproducible taxonomy to avoid counting relabeling as invention. For a causal separation of displacement and complementarity, specify distinct counterfactual technology components and an interaction-allocation rule; aggregate employment signs alone cannot identify the mechanisms. Distributional and welfare conclusions additionally require ownership, price, and household-preference assumptions.
\subsection*{Source audit and status}
[S1] concerns earlier computerization, [S2] industrial robots, and [S3] a workplace generative-AI deployment. The atlas's 2023 citation is NBER Working Paper 31161; the journal article appeared in 2025. Existing local productivity and robot results must not be relabeled unsolved. The dynamic, cross-setting equilibrium prediction problem is editorial and conditional on $M$, with no claim that every instance remains open in 2026.

\subsection*{References}
\begin{enumerate}[label=\textnormal{[S\arabic*]},leftmargin=*]
\item David H. Autor, Frank Levy, and Richard J. Murnane (2003). \emph{The Skill Content of Recent Technological Change: An Empirical Exploration}. The Quarterly Journal of Economics 118(4), 1279--1333. \url{https://doi.org/10.1162/003355303322552801}. \textit{Locator:} Primary publisher, working-paper, or author record: bibliographic information and abstract; full-text theorem verification is not claimed. \textit{Role:} Task-based empirical evidence; background for task substitution.
\item Daron Acemoglu and Pascual Restrepo (2020). \emph{Robots and Jobs: Evidence from US Labor Markets}. Journal of Political Economy 128(6), 2188--2244. \url{https://doi.org/10.1086/705716}. \textit{Locator:} Primary publisher, working-paper, or author record: bibliographic information and abstract; full-text theorem verification is not claimed. \textit{Role:} Local labor-market evidence for industrial robots; cannot by itself identify generative-AI equilibrium effects.
\item Erik Brynjolfsson, Danielle Li, and Lindsey R. Raymond (2023). \emph{Generative AI at Work}. NBER Working Paper 31161; published as Generative AI at Work, The Quarterly Journal of Economics 140(2), 889--942 (2025). \url{https://www.nber.org/papers/w31161}. \textit{Locator:} Primary publisher, working-paper, or author record: bibliographic information and abstract; full-text theorem verification is not claimed. \textit{Role:} Workplace AI deployment evidence; a local productivity result rather than an aggregate wage theorem. The 2023 date denotes the working paper, not the 2025 journal article; publisher DOI 10.1093/qje/qjae044 verified in turn76search9.
\end{enumerate}

\subsection*{Common conventions: Source mathematical conventions}

Unless an entry states otherwise, agent and firm sets are finite. State and action spaces are measurable, strategies and outcome maps are measurable, probability laws are specified, and expectations exist for the targets under discussion. Infinite-horizon discount factors lie in $(0,1)$ and discounted objectives are finite. Norms, tolerances and complexity input sizes must be specified for computational comparisons. Constants or primitives that appear locally are inputs, not empirical facts asserted by this catalogue.

\textbf{Identification.} A statistical model $m\in\mathcal M$ implies an observable law $P_m$ and target $\theta(m)$. At law $P$, its identified set is
\[
\Theta(P;\mathcal M)=\{\theta(m):m\in\mathcal M,\ P_m=P\}.
\]
Point identification holds when this set is a singleton, or equivalently when $P_{m_1}=P_{m_2}$ implies $\theta(m_1)=\theta(m_2)$ for all compatible models. Sharp bounds mean bounds attained, or approached when the boundary is not attained, by compatible models. Writing this set is a definition, not a solution: the substantive task is to establish informative restrictions and derive or estimate the set. An unrestricted-confounding model may give only trivial bounds.

\textbf{Inference.} Identification, estimability and valid uncertainty quantification are separate questions. Fix $\alpha\in(0,1)$. A confidence procedure $C_n$ over a declared law class $\mathcal P$ has asymptotic uniform coverage at level $1-\alpha$ when
\[
\liminf_{n\to\infty}\inf_{P\in\mathcal P}\Pr_P\{\theta(P)\in C_n\}\geq1-\alpha.
\]
For set-identified targets the coverage event must be defined separately, for example coverage of the full identified set. If an entry uses identification-indexed models instead of uniquely indexed laws, the same coverage convention is applied over those models. Pointwise limits do not establish uniform validity, and asymptotic validity does not guarantee finite-sample precision. Sample sizes, dependence structures and weak-identification sequences are part of each application.

\textbf{Counterfactuals.} $Y_i(a)$ is a potential outcome under a specified intervention, including its timing, implementation and versions. It is not $Y_i$ conditional on voluntarily choosing $a$. A full intervention vector permits interference; shorthand scalar treatment notation does not silently impose no interference. Assignment, selection, attrition, anticipation and measurement must be specified. A claim about an instrument's local effect is not automatically a population policy effect.

\textbf{Economic equilibrium.} A counterfactual model includes best responses, constraints, feasibility, market clearing, state transitions and expectations consistent with the stated information. When several equilibria exist, either a declared selector is an input or the target is set-valued. Comparisons across policy regimes must use the stated selection convention. Reduced-form relationships are not presumed invariant after an equilibrium-changing intervention.

\textbf{Optimization and welfare.} Optimization is over the stated feasible set. If attainment is not established, $\sup$ or $\inf$ and $\varepsilon$-optimal choices replace an asserted maximizer or minimizer. For example, with value $V=\sup_{a\in\mathcal A}W(a)<\infty$, an $\varepsilon$-optimal set is $\{a\in\mathcal A:W(a)\geq V-\varepsilon\}$ for $\varepsilon>0$. Benefits, costs, risk and health or environmental outcomes can be aggregated only using declared units and conversion weights. Interpersonal utility weights, the treatment of fiscal transfers and foreign profits, discounting, and the behavioral welfare criterion are explicit inputs. A comparative-static derivative additionally requires existence and appropriate differentiability of the selected counterfactual.

\textbf{Sources.} Local labels [S1], [S2], and so forth restart in every question. Locators identify the relevant abstract, model, section or result at the level actually checked; a bibliographic or abstract-level verification is not represented as inspection of an entire proof. Working papers are distinguished from later publications and changing versions. Source summaries are paraphrased. All URL checks and editorial formulations refer to the audit date stated above.
% END ECONOMICS STATEMENT
\end{document}
